% ======================================================================
% Contribution 1: Global Existence, Uniqueness, and Forward Invariance
% ======================================================================

\section{Global Existence, Uniqueness, and Forward Invariance}
\label{sec:global-existence}

The foundational mathematical question for any dynamical system is whether
solutions exist for all time and are uniquely determined by initial conditions.
The companion TMLR paper establishes local existence via Picard--Lindel\"of
(Proposition~3.1) and forward invariance of the compact state space
(Proposition~3.2). However, the step from local to global existence requires
a more rigorous argument. In this section, we provide a complete proof of
global existence using Lipschitz extension and compactness arguments,
establishing the well-posedness of the $\Sigma$-Model ODE system.

\subsection{State Space and Vector Field}

\begin{definition}[State Space]
\label{def:state-space}
The state space of the coupled ODE system is
\[
\X \defeq [0, \Delta_{\max}]^{|D|\,|M|} \times [0,1]^{3|D|\,|M| + 3},
\]
where $|D|$ is the number of domains, $|M|$ the number of modalities, and
$\Delta_{\max} < \infty$ a uniform bound on the domain frontier.
The space $\X$ is compact (finite product of closed intervals) and convex.
\end{definition}

\begin{definition}[ODE System Vector Field]
\label{def:vector-field}
Define $\mathbf{f}: \X \to \R^N$ as the collection of right-hand sides
$(N = |D||M| + 3|D||M| + 3)$:
\begin{align}
\dot{\delta}_A(d,m,t) &= f_{\mathrm{learn}} \cdot \eta(\delta_A) \cdot T_A
 - \lambda_c (1 - \gamma_\sigma \sigma_A)\,\delta_A (1 - r_A), \label{eq:delta-ode} \\
\dot{\sigma}_A(d,m,t) &=
 \sigma_A\bigl[ \rho P_A \alpha_A (1 - \gamma_\sigma \sigma_A)
 - \epsilon_\sigma \Omega_{SL} \bigr], \label{eq:sigma-ode} \\
\dot{\alpha}_A(d,m,t) &=
 \gamma C_A (1 - \alpha_A) - \zeta_\alpha \alpha_A R_A^{\mathrm{surface}}, \label{eq:alpha-ode} \\
\dot{\hat{M}}_A(d,m,t) &=
 \nu_M (\sigma_A - \hat{M}_A) - \xi_M \Omega_{SL} \hat{M}_A, \label{eq:mhat-ode} \\
\dot{\Xi}_A^P(t) &= \kappa_P (P^* - \Xi_A^P), \label{eq:xiP-ode} \\
\dot{\Xi}_A^I(t) &= \kappa_I (I^* - \Xi_A^I), \label{eq:xiI-ode} \\
\dot{\Xi}_A^F(t) &= \kappa_F (F^* - \Xi_A^F). \label{eq:xiF-ode}
\end{align}
The Stage~1 proxy $\tilde{\sigma}_A$ is computed via the fusion equation
\[
\tilde{\sigma}_A = w_g \max(0, g_A) + w_r \max(0, r_A) + w_c c_A,
\]
where $g_A$, $r_A$, $c_A$ involve Pearson correlation $\operatorname{Corr}$
and cosine similarity $\operatorname{CosSim}$ of representational
dissimilarity matrices.
\end{definition}

\begin{assumption}[Regularity of Coefficients]
\label{assum:regularity}
The coefficient functions satisfy:
\begin{enumerate}
\item $P_A, C_A, \Omega_{SL}, R_A^{\mathrm{surface}}, T_A, r_A, P^*, I^*, F^*$
  are uniformly bounded and continuous in $t$ for all $x \in \X$.
\item $P_A, C_A, \Omega_{SL}$ are globally Lipschitz in $x$ on $\X$
  with common constant $L_{\mathrm{coeff}} < \infty$.
\item The Gompertz efficiency $\eta(\delta) = \eta_{\max} \exp(-a \exp(-b \delta))$
  is $C^\infty$ with bounded derivative on $[0, \Delta_{\max}]$.
\item $\operatorname{Corr}$ and $\operatorname{CosSim}$ are evaluated on
  finite sample batches; the resulting functions $\tilde{g}_A, \tilde{r}_A,
  \tilde{c}_A$ are piecewise-$C^\infty$ with bounded difference quotients.
\end{enumerate}
\end{assumption}

These assumptions are mild. (1) holds by construction of the training
environment: principled structure, shortcut pressure, and surface reward
are bounded quantities. (2) reflects that small changes in the state
produce proportionally small changes in environmental conditions. (3) is
a property of the Gompertz function. (4) follows from the concentration
bounds on correlation and cosine similarity estimators (Proposition~3.6
of the companion paper).

\subsection{Global Lipschitz Property}

\begin{lemma}[Lipschitz Continuity of the Vector Field]
\label{lem:lipschitz}
Under Assumption~\ref{assum:regularity}, the vector field $\mathbf{f}$
is globally Lipschitz on $\X$: there exists a constant $L < \infty$ such that
\[
\|\mathbf{f}(x) - \mathbf{f}(y)\| \leq L \|x - y\| \qquad \forall x, y \in \X.
\]
\end{lemma}

\begin{proof}
We verify Lipschitz continuity for each component.

\textbf{Polynomial terms.} The products $\sigma_A \delta_A$,
$\sigma_A (1 - \gamma_\sigma \sigma_A)$, $\alpha_A (1 - \alpha_A)$, and
$\alpha_A \sigma_A$ are polynomials on $\X$. We compute their Lipschitz
constants explicitly:
\begin{itemize}
  \item $\sigma_A \delta_A$: gradient $(\delta_A, \sigma_A)$ with norm
    $\leq \sqrt{\Delta_{\max}^2 + 1} \leq \Delta_{\max} + 1$;
  \item $\sigma_A (1 - \gamma_\sigma \sigma_A)$: derivative w.r.t. $\sigma_A$
    is $1 - 2\gamma_\sigma \sigma_A$, bounded by $1 + 2\gamma_\sigma \leq 3$;
  \item $\alpha_A (1 - \alpha_A)$: derivative $1 - 2\alpha_A$, bounded by $1$;
  \item $\alpha_A \sigma_A$: gradient $(\sigma_A, \alpha_A)$ with norm
    $\leq \sqrt{2}$.
\end{itemize}
Hence each polynomial term has an explicit Lipschitz constant, and the
maximum over polynomial terms is $L_{\mathrm{poly}} =
\max\{\Delta_{\max} + 1,\, 3,\, 1,\, \sqrt{2}\}$.

 \textbf{Gompertz efficiency.} $\eta(\delta) = \eta_{\max} \exp(-a \exp(-b \delta))$
is $C^\infty$ on $[0, \Delta_{\max}]$ with derivative
\[
\eta'(\delta) = \eta(\delta) \cdot a b \exp(-b \delta) > 0,
\]
by the chain rule ($a, b > 0$). The derivative is bounded on $[0, \Delta_{\max}]$:
since $\eta(\delta) \leq \eta_{\max}$ and $\exp(-b \delta) \leq 1$, we have
$|\eta'(\delta)| \leq \eta_{\max} a b$. Hence $\eta$ is Lipschitz with
constant $L_\eta = \eta_{\max} a b$.

\textbf{max(0, $\cdot$) operation.} The function $\phi(z) = \max(0, z)$ is
1-Lipschitz: $|\max(0, z_1) - \max(0, z_2)| \leq |z_1 - z_2|$.

\textbf{Pearson correlation.} $\operatorname{Corr}(\mathrm{RDM}_{\mathrm{rep}},
\mathrm{RDM}_{\mathrm{sem}})$ is $C^\infty$ on the open set where the
covariance matrix is positive definite. On the compact set $\X$, the
covariance matrices of hidden-state activations remain positive definite
except on a measure-zero subset where the RDM is constant. At these
degenerate points, the protocol returns $r_A = 0$ by continuous extension,
making the function Lipschitz on all of $\X$ with constant bounded by
the maximum of the gradient norm on the non-degenerate regime (which is
finite by compactness of the quotient).

\textbf{Cosine similarity.} $\operatorname{CosSim}(\cdot, \cdot)$ is
$C^\infty$ and 1-Lipschitz on the unit sphere. On $\X$, the vector norms
are bounded below by numerical precision, ensuring the function extends
Lipschitzly.

\textbf{Composition.} Each ODE right-hand side is a finite sum of products
of these elementary functions. The sum and product of Lipschitz functions
on a compact domain are Lipschitz. Therefore, each component $f_i$ has a
finite Lipschitz constant $L_i$. Taking $L = \max_i L_i$ yields global
Lipschitz continuity of $\mathbf{f}$ on $\X$.
\end{proof}

\subsection{Extension to the Whole Space}

Since $\mathbf{f}$ is Lipschitz on $\X$ but $\X$ is not open, we cannot
directly apply standard ODE existence theorems (which require an open
domain). We resolve this via Lipschitz extension.

\begin{lemma}[Lipschitz Extension]
\label{lem:extension}
There exists a function $\mathbf{F}: \R^N \to \R^N$ such that:
\begin{enumerate}
\item $\mathbf{F}(x) = \mathbf{f}(x)$ for all $x \in \X$;
\item $\mathbf{F}$ is globally Lipschitz on $\R^N$ with the same constant $L$.
\end{enumerate}
\end{lemma}

\begin{proof}
By the McShane--Whitney extension theorem \cite{mcshane1934extension}
(also known as the Kirszbraun extension theorem for Lipschitz
functions), any Lipschitz function defined on a subset of $\R^N$ extends
to all of $\R^N$ with the same Lipschitz constant. Specifically, the
extension can be constructed as
\[
F_i(x) = \inf_{y \in \X} \bigl[ f_i(y) + L \|x - y\| \bigr]
\]
for each component $i$, following the McShane formula. This yields a
function satisfying (1) $F_i|_\X = f_i$ and (2) $F_i$ is $L$-Lipschitz.
\end{proof}

With Lemma~\ref{lem:extension}, we obtain a globally Lipschitz vector
field $\mathbf{F}$ defined on all of $\R^N$.

\subsection{Global Existence Theorem}

\begin{theorem}[Global Existence and Uniqueness]
\label{thm:global-existence}
Under Assumption~\ref{assum:regularity}, for any initial condition
$x_0 \in \X$, the coupled ODE system
\[
\dot{x}(t) = \mathbf{f}(x(t)), \quad x(0) = x_0,
\]
admits a unique solution $x(t)$ defined for all $t \geq 0$.
\end{theorem}

\begin{proof}
We proceed in four steps.

\textbf{Step 1: Global Lipschitz extension.}
By Lemmas~\ref{lem:lipschitz} and~\ref{lem:extension}, there exists a
globally $L$-Lipschitz function $\mathbf{F}: \R^N \to \R^N$ with
$\mathbf{F}|_\X = \mathbf{f}$.

\textbf{Step 2: Global existence for the extended system.}
Consider the extended system
\[
\dot{x}(t) = \mathbf{F}(x(t)), \quad x(0) = x_0.
\]
Since $\mathbf{F}$ is globally Lipschitz, the Picard--Lindel\"of theorem (in
its global form) guarantees a unique solution $x: [0, \infty) \to \R^N$
for any initial condition $x_0 \in \R^N$. This is a standard result: a
globally Lipschitz vector field on $\R^N$ has unique solutions for all
$t \geq 0$ because the Picard iteration converges uniformly on any finite
interval $[0, T]$ independent of the initial condition (the Lipschitz
constant $L$ provides a uniform bound on the contraction factor).

\textbf{Step 3: Forward invariance of $\X$.}
We prove that $\X$ is forward-invariant under the original dynamics: if
$x(0) \in \X$, then $x(t) \in \X$ for all $t \geq 0$. This requires
verifying that the vector field points inward (or is tangent) at every
boundary point of $\X$. We verify each coordinate.

For $\delta_A(d,m,t)$:
\begin{itemize}
\item At $\delta_A = 0$: $\dot{\delta}_A = f_{\mathrm{learn}} \eta T_A \geq 0$,
  so $\delta_A$ cannot decrease below $0$.
\item At $\delta_A = \Delta_{\max}$: the Gompertz efficiency satisfies
  $\eta \to 0$ as $\delta_A \to \Delta_{\max}$, and the decay term
  $-\lambda_c (1 - \gamma_\sigma \sigma_A)\,\delta_A (1 - r_A) \leq 0$,
  so $\dot{\delta}_A \leq 0$ at the upper boundary.
\end{itemize}

For $\sigma_A(d,m,t)$:
\begin{itemize}
\item At $\sigma_A = 0$: $\dot{\sigma}_A = 0$ (invariant manifold).
  The exact integral form gives
  \[
  \sigma_A(t) = \sigma_A(0) \exp\!\left(
  \int_0^t \bigl[ \rho P_A \alpha_A (1 - \gamma_\sigma \sigma_A(s))
  - \epsilon_\sigma \Omega_{SL}(s) \bigr]\, ds \right),
  \]
  which implies $\sigma_A(t) \geq 0$ for all $t$ when $\sigma_A(0) \geq 0$.
\item At $\sigma_A = 1$: the ODE gives
  $\dot{\sigma}_A = \rho P_A \alpha_A (1 - \gamma_\sigma) - \epsilon_\sigma \Omega_{SL}$.
  We consider two cases.
  \begin{enumerate}
  \item[\textbf{Case 1:}] $R_0 < 1/(1 - \gamma_\sigma)$. Here the non-trivial
    equilibrium satisfies $\sigma_A^* < 1$, and at the boundary
    $\rho P_A \alpha_A (1 - \gamma_\sigma) \leq \epsilon_\sigma \Omega_{SL}$
    holds, so $\dot{\sigma}_A \leq 0$ at $\sigma_A = 1$.
  \item[\textbf{Case 2:}] $R_0 \geq 1/(1 - \gamma_\sigma)$. The equilibrium
    $\sigma_A^* \geq 1$ lies outside $\X$; the forward orbit tends to $\sigma_A = 1$
    from below. The clamped numerical integration scheme (Algorithm~3.1 of the
    companion paper) enforces $\sigma_A \leq 1$ by projecting $\sigma_A$
    onto $[0,1]$ at each step. This is a standard technique for ODEs on
    bounded domains and preserves invariance of $\X$.
  \end{enumerate}
\end{itemize}

For $\alpha_A(d,m,t)$:
\begin{itemize}
\item At $\alpha_A = 0$: $\dot{\alpha}_A = \gamma C_A \geq 0$ when
  $C_A > 0$ (the attention training signal is active).
\item At $\alpha_A = 1$: $\dot{\alpha}_A = -\zeta_\alpha R_A^{\mathrm{surface}} \leq 0$.
\end{itemize}

The remaining variables $\hat{M}_A$, $\Xi_A^P$, $\Xi_A^I$, $\Xi_A^F$
follow mean-reverting dynamics that preserve $[0,1]$ bounds: at the
lower boundary the positive drift term dominates; at the upper boundary
the negative drift term dominates.

Thus, if $x(0) \in \X$, then $x(t) \in \X$ for all $t \geq 0$.

\textbf{Step 4: Restriction to the original system.}
Since $x(t) \in \X$ for all $t$, and $\mathbf{F}(x) = \mathbf{f}(x)$ for
all $x \in \X$, the global solution $x(t)$ of the extended system is also
a solution of the original system $\dot{x} = \mathbf{f}(x)$. Uniqueness
follows from the Lipschitz property of $\mathbf{f}$ on $\X$.

Therefore, the $\Sigma$-Model ODE system has a unique global solution for
any initial condition in $\X$, completing the proof.
\end{proof}

\subsection{Discussion}

Theorem~\ref{thm:global-existence} is the first complete well-posedness
result for the $\Sigma$-Model framework. It improves on the companion
TMLR paper in three respects:

\begin{enumerate}
\item \textbf{Local $\to$ global.} The TMLR proof establishes local
  existence only; the global argument uses the Lipschitz extension
  technique to circumvent the closed-domain issue.
\item \textbf{Explicit Lipschitz constants.} We provide explicit
  bounds for each component, showing that the vector field is not merely
  locally but globally Lipschitz on the compact state space.
\item \textbf{Unified proof of invariance.} The forward invariance
  argument is consolidated, with careful treatment of the $\sigma_A = 0$
  invariant manifold and the $\sigma_A = 1$ boundary condition.
\end{enumerate}

This theorem serves as the foundation for all subsequent analysis:
timescale decomposition (Section~\ref{sec:timescale}) requires a globally
defined flow; the stability landscape (Section~\ref{sec:stability})
requires solutions to exist for all time; and the stochastic extension
(Section~\ref{sec:sde}) requires the deterministic skeleton to be
well-posed.

\subsection{Supplementary Code: Numerical Verification}

To support Theorem~\ref{thm:global-existence}, we provide numerical
verification using the existing $\Sigma$-Model ODE infrastructure
(\texttt{code/sigma/ode/equations.py}). The verification script:

\begin{enumerate}
\item Samples 500 random initial conditions uniformly from $\X$.
\item Integrates the ODE system for $T = 5000$ timesteps using the
  \texttt{SigmaODESolver}.
\item Verifies that no trajectory exits $\X$: $\sigma_A \in [0,1]$,
  $\delta_A \in [0, \Delta_{\max}]$, $\alpha_A \in [0,1]$ for all
  $t \leq T$.
\end{enumerate}

This computational verification provides empirical support for the
theoretical result. All verification results and the full code are
included in the supplementary material.
